\section{Introduction and statement of results}\label{sec:intro}

A classical approach to obtain and prove $q$-series identities is the Bailey lemma, originally found by Bailey~\cite{Ba}, and whose iterative strength was later highlighted by Andrews~\cite{A1, A2, AAR} through the so-called Bailey chain.
Fix complex numbers $a$ and $q$. Recall~\cite{Ba} that a Bailey pair $((\alpha_n)_{n\geq0}, (\beta_n)_{n\geq0})$ ($(\alpha_n, \beta_n)$ for short) relative to $a$ is a pair of sequences satisfying
\begin{equation}\label{bp}
\beta_n=\sum_{j=0}^n\frac{\alpha_j}{(q)_{n-j}(aq)_{n+j}}\qquad\forall n\in\mathbb{N}.
\end{equation}
Here and throughout the paper, we use standard $q$-series notations which can be found in~\cite{GR}
\begin{equation*}
(a)_\infty = (a;q)_\infty:=\prod_{j\geq 0}\bigl(1-aq^j\bigr)\qquad\mbox{and}\qquad(a)_k = (a;q)_k:=\frac{(a;q)_\infty}{\bigl(aq^k;q\bigr)_\infty},
\end{equation*}
where $k \in \Z$, and
$
(a_1,\ldots,a_m)_k:=(a_1)_k\cdots(a_m)_k$,
where $k$ is an integer or infinity, and as usual~${|q|<1}$ to ensure convergence of infinite products.
